\subsection{上中心列}

设 \(\mathfrak{g}\) 是李代数，\(\mathbf{P}\) 是 \(\mathfrak{g}\) 的子集。\(\mathbf{P}\) 在 \(\mathfrak{g}\) 中的中心化子是由 \(\mathfrak{g}\) 中与 \(\mathbf{P}\) 中元素可交换的元素组成的集合。此中心化子是所有 ad \(y\) 的核的交，其中 \(y\) 遍历 \(\mathbf{P}\)；因此它是 \(\mathfrak{g}\) 的子代数。

命题 5. 设 g 是李代数，a 是 g 的理想（分别为特征理想）。a 在 g 中的中心化子 \(\mathfrak{a}'\) 是 g 的理想（分别为特征理想）。

设 \(D\) 是 \(g\) 的内导子（分别为导子）。若 \(x \in a'\) 且 \(y \in a\)，则

\[
[ \mathrm{D} x, y ] = \mathrm{D} ([ x, y ]) - [ x, \mathrm{D} y ] = 0;
\]

因此 \(\mathrm{D}x\in \mathfrak{a}^{\prime}\)。故命题得证。

设 g 是李代数。g 在自身中的中心化子称为 g 的中心，即由 \(x \in g\) 满足 \([x, y] = 0\) 且对所有 \(y \in g\) 成立的元素构成的特征理想。g 的中心是同态 \(x \mapsto ad x\) 的核。

\(\mathfrak{g}\) 的上中心列是递增序列 \(\mathcal{C}_{0}\mathfrak{g},\mathcal{C}_{1}\mathfrak{g},\ldots\)，其各项是 \(\mathfrak{g}\) 的特征理想，递归定义如下：(1) \(\mathcal{C}_{0}\mathfrak{g} = \{0\}\)；(2) \(\mathcal{C}_{p + 1}\mathfrak{g}\) 是在从 \(\mathfrak{g}\) 到 \(\mathfrak{g} / \mathcal{C}_{p}\mathfrak{g}\) 的典范映射下，\(\mathfrak{g} / \mathcal{C}_{p}\mathfrak{g}\) 的中心的逆像。

理想 \(\mathcal{C}_{1}\mathfrak{g}\) 是 \(\mathfrak{g}\) 的中心。

